% Sixth Philippic --- headnote
% philippic-6, 4 January 43 BC, Rome

\textit{The popular counterpart to the Fifth Philippic, delivered
to the people in the Forum on 4 January 43 BC, three days after
the great senatorial debate of the Kalends. The Senate had met
on 1 January under the new consuls Aulus Hirtius and Gaius
Vibius Pansa, in the temple of Concord, to take up the war
against Antony. Cicero (Philippic 5) had argued for an immediate
declaration of a \textit{tumultus}, with \textit{iustitium} and
\textit{saga} and a levy --- the full apparatus of war ---
against Antony at Mutina. Quintus Fufius Calenus, the
consular father-in-law of Pansa, had moved instead that
legates be sent. The debate ran for three days; today, on
4 January, the Senate had at last carried the legation, sending
Servius Sulpicius Rufus, Lucius Calpurnius Piso, and Lucius
Marcius Philippus to Antony with the terms Antony was to obey:
quit Mutina, quit Gaul, withdraw to within two hundred miles of
Rome, submit himself to the will of the Senate and the Roman
people. Cicero had lost his motion. The Sixth Philippic is the
speech he gave the same day, in the Forum, to report to the
people what had been decided and to explain why they should
nevertheless not lose heart.}

\textit{The shape is short --- nineteen sections --- and divided
into three unequal parts. \textsection 1--9 are the report and
the case against the legation: the matter has been decided, but
\textit{less severely than was fitting}, the war is delayed but
not its cause removed (\textsection 1); the legation is in
substance a notice of war, like the embassy the Senate once
sent to Hannibal not to attack Saguntum (\textsection 4, 6);
Antony will not obey, for he has never obeyed his own judgment,
let alone anyone else's (\textsection 4--5); the brigand-
gladiator language of the Fourth Philippic is renewed
(\textit{gladiator}, \textsection 3; \textit{importunissima
belua}, \textsection 7); Decimus Brutus, the consul-elect now
penned in Mutina, is the test --- a Brutus born for liberty and
not for his own leisure (\textsection 8--9). \textsection 10--15
turn aside into the catalogue of Lucius Antonius's statues ---
the murmillo brother who is the \textit{Africanus} of the
Antonine inner circle --- a sustained set-piece of public
ridicule: the gilded equestrian statue erected by the thirty-
five tribes inscribed \textit{To their patron from the thirty-
five tribes} (\textsection 12); the Forum statue, in front of
the temple of Castor, beside that of Quintus Tremulus who
conquered the Hernici (\textsection 13); the equestrian-order
statue with its inscription \textit{to their patron}
(\textsection 13); the military-tribunes' statue, set up by men
who had served twice in Caesar's army and to whom Antony had
divided up the Semurium estate (\textsection 14); and the
\textit{palmaris} statue, the consummate one, by the
moneylenders of the Middle Janus, inscribed \textit{To Lucius
Antonius, from the Middle Janus, their patron} (\textsection
15). Each is named by name, each ridiculed. The Trebellius
section (\textsection 11) is the bridge: a man who took
\textit{fides} as his cognomen and lives by defrauding his
creditors. \textsection 16--19 close with the hortatory turn:
wait a few days for the legates to come back; when they bring
back the inevitable report of Antony's refusal, the case for
war will be made for us; meanwhile the unanimity of the people,
the Senate, the orders, the colonies, all Italy, is so great
that no hour more can be delayed. The final word is
\textit{libertas}: \textit{Other nations can bear slavery;
liberty belongs to the Roman people by right.}}

\textit{The register is the register of the contio: shorter
periods than the senatorial debate, named persons named, jokes
told (the Trebellius \textit{O fides!}, the imagined moneylender
of the Middle Janus, the Mylasa murmillo cutting his friend's
throat), the structural drumbeat of \textit{Quirites} sustained
through almost every section. The speech runs throughout on the
tension between what Cicero proposed in the Senate and what the
Senate carried: \textit{minus severe quam decuit, non tamen
omnino dissolute} (\textsection 1) is the controlled, slightly
bitter judgment that opens the speech and is restated at
\textsection 7 (\textit{non omnino dissolutum est quod decrevit
senatus: habet atrocitatis aliquid legatio: utinam nihil
haberet morae}). Cicero is making the best of a compromise he
opposed; the report to the people is also a holding action ---
a few days only, the cloaks will be put on, the war will come,
and meanwhile the man being notified is precisely the man who
will not be notified.}

\textit{The legation in fact returned almost at once. Servius
Sulpicius Rufus, the great jurist, died on the road; Piso and
Philippus reached Antony and were rebuffed --- Antony added his
own counter-demands (consular province of Further Gaul, the
veteran lands, immunity for his acts) and refused the Senate's
terms. By 2 February the legates were back; on 3 February the
Senate, addressed by Cicero in the Eighth Philippic, finally
decreed \textit{tumultus}. The drum-rhythm Cicero set going in
this contio --- \textit{a few days, then the cloaks} ---
played out exactly as he had told the people it would.}
